% ECONOMICS_PROBLEM: {"id": "F16-250", "title": "Distribution of trade gains", "jel_code": "F16", "source_id": "OP-0340"}
% !TeX program = xelatex
\documentclass[11pt,a4paper]{article}
\usepackage[margin=23mm]{geometry}
\usepackage{fontspec}
\setmainfont{Latin Modern Roman}
\usepackage{amsmath,amssymb,mathtools}
\usepackage{xcolor,enumitem,booktabs,longtable,array,xurl,hyperref}
\definecolor{muted}{HTML}{53636C}
\hypersetup{colorlinks=true,urlcolor=blue,linkcolor=blue}
\setlength{\parindent}{0pt}
\setlength{\parskip}{5pt}
\newcommand{\R}{\mathbb R}
\newcommand{\E}{\mathbb E}
\newcommand{\N}{\mathbb N}
\newcommand{\ind}{\mathbf 1}
\newcommand{\Var}{\operatorname{Var}}
\newcommand{\Cov}{\operatorname{Cov}}
\newcommand{\supp}{\operatorname{supp}}
\DeclareMathOperator*{\argmin}{arg\,min}
\DeclareMathOperator*{\argmax}{arg\,max}
\newcommand{\entrymeta}[3]{{\sffamily\small\color{muted}\textbf{\texttt{#1}}\quad|\quad #2\par #3\par}}
\newcommand{\Problemref}[1]{\texttt{#1}}
\newcommand{\JELref}[2]{\texttt{#2} (JEL \texttt{#1})}
\begin{document}
\section*{Distribution of trade gains}
\textbf{Problem ID:} \texttt{F16-250}\quad \textbf{JEL:} \texttt{F16}\par
% BEGIN ECONOMICS STATEMENT
\label{problem:OP-0340}
{\sffamily\small\color{muted}Original subject group: Trade, globalization and geoeconomics\par}
\label{definitions:problem:30007622}
\textbf{Audit classification:} Background; proposed extension. The citation mixed the 2021 NBER WP with a 2024 journal URL. Both are real but distinct versions and are now separately listed. The exact household EV distribution problem is editorial, not verified as a canonical open question.
\subsection*{Definitions and precise research target}
Fix households \(i\), industries, locations and a trade-cost intervention \(\tau\) relative to baseline \(0\). Let \(U_i(\tau)\) be equilibrium lifetime utility under wages, employment opportunities, consumer prices and wealth returns induced by the intervention. Define equivalent variation \(EV_i(\tau)\) as the lump-sum amount added to baseline wealth that equates baseline attainable utility to \(U_i(\tau)\), using the declared lifetime budget and preference model. The task is to identify the distribution
\[F_{\tau}(v)=\Pr(EV_i(\tau)\leq v)\]
for a stated household population, and decompose changes by earnings, expenditures and asset-ownership channels under explicitly ordered counterfactuals. Include migration, occupational switching and general-equilibrium price responses when the horizon permits them. Specify household preference heterogeneity and the assumptions translating measured income and price changes into welfare. Establish the exogeneity of the trade shock or derive an identified set when trade policy responds to domestic conditions. Report uncertainty and dependence on channel ordering, since interactions generally prevent a unique additive decomposition. An aggregate gains estimate or a local wage effect does not identify the full joint welfare distribution. Separate transitional losses from permanent effects and disclose exclusions caused by missing household wealth data.

\textbf{Additional formal definitions and scope (editorial).} Let \(v_i(p^0,w_i)\) be baseline lifetime indirect utility at the entire baseline price path \(p^0\), with initial wealth \(w_i\). Equivalent variation solves \(v_i(p^0,w_i+EV_i)=U_i(\tau)\); assume strict monotonicity in wealth and that the target is in its range. Positive \(EV_i\) denotes a gain. Compensating variation instead changes wealth at the postintervention prices and is generally unequal. Fix a probability measure \(\mu\) over baseline households, including the treatment of births, deaths and migrants, so \(F_\tau(v)=\mu\{i:EV_i\leq v\}\). Identify the joint law of household preferences, incomes, spending and assets before constructing this distribution: marginal income or price distributions alone do not determine it. Counterfactual channel changes must retain budget feasibility and use an explicit order.

For the identification part, declare an admissible class \(\Theta\) of structural models, including preferences, technologies, shock laws, measurement/sampling rules and any equilibrium selector. If \(O\) denotes the complete observable history and \(T(\theta)\) the target defined above, the identified set is \(\mathcal I_T=\{T(\theta):\theta\in\Theta,\ \mathcal L_\theta(O)=\mathcal L_{\rm obs}(O)\}\); point identification means this set is a singleton. A \(do(a)\) comparison replaces stated policy/technology equations while fixing the declared exogenous law and re-solving the remaining equations. These definitions do not assert that an identification theorem holds without further restrictions.
\subsection*{Variants and assumption changes}
\begin{enumerate}
\item \textbf{Static goods-price and wage effects (source-motivated).} Fix households and occupations and recover equivalent variation from observed preferences and joint income/expenditure changes.
\item \textbf{Lifetime adjustment and wealth (editorial).} Allow moving, occupational investment and asset returns; identify the household welfare distribution and ordered channel contributions.
\end{enumerate}
\subsection*{References and source audit}
\begin{enumerate}
\item NBER verifies Rodrik, 2021 WP 29507 and DOI; it is distinct from the 2024 journal publication originally linked. Locator: 2021 NBER abstract; author-hosted version, compensation discussion printed pp.3--4. Support: Trade-distribution and compensation background. NBER indexed metadata and author-hosted text inspected. URL: \url{https://www.nber.org/papers/w29507}.
\item Publisher verifies Rodrik, 2024, Oxford Open Economics 3(Supplement 1):i1076--i1082, DOI 10.1093/ooec/odad048; not the 2021 NBER DOI. Locator: Compensation discussion; author-hosted version printed pp.3--4. Support: Compensation background; no exact optimal-program open theorem. Publisher indexed text and author-hosted version inspected; repeated publisher fetch failed. URL: \url{https://academic.oup.com/ooec/article/3/Supplement_1/i1076/7708093}.
\end{enumerate}
\begin{enumerate}\raggedright
\item\hypertarget{definitions:source:30007622:1}{} Dani Rodrik. 2021. \emph{A Primer on Trade and Inequality}. 2021 NBER abstract; author-hosted version, compensation discussion printed pp.3--4.
\url{https://www.nber.org/papers/w29507}.
DOI: \url{https://doi.org/10.3386/w29507}.
\item\hypertarget{definitions:source:30007622:2}{} Dani Rodrik. 2024. \emph{A primer on trade and inequality}. Compensation discussion; author-hosted version printed pp.3--4.
\url{https://academic.oup.com/ooec/article/3/Supplement_1/i1076/7708093}.
DOI: \url{https://doi.org/10.1093/ooec/odad048}.
\end{enumerate}

\subsection*{Common conventions: Source mathematical conventions}

These conventions apply to each entry unless explicitly replaced.

\textbf{Models and identification.} A primitive model \(m\) specifies all probability laws, feasible choices, information, equilibrium and measurement rules used by its target. The admissible class \(\mathcal M\) is fixed independently of outcomes. If \(P_O(m)\) is its observable probability law and \(\tau(m)\) the target, its identified set at observable law \(P\) is
\[
 \mathcal I_\tau(P)=\{\tau(m):m\in\mathcal M,\ P_O(m)=P\}.
\]
Point identification means that this set is a singleton. An empty set signals incompatibility of data and assumptions. Coordinate bounds need not describe the joint set. Bounds are sharp only if every retained target is attainable or arbitrarily approachable by compatible models. Finite bounds require explicit restrictions on unobserved quantities; unrestricted confounding, hidden wealth, extrapolation or arbitrary equilibrium selection can yield uninformative sets.

\textbf{Causal interventions.} A potential outcome \(Y_i(p)\) is the outcome under a complete policy regime \(p\), including timing, financing, information and market spillovers. The target population and units are fixed before treatment. With interference, use the complete assignment vector or a justified exposure mapping; individual no-interference notation alone cannot identify an economy-wide intervention. A difference of mean potential outcomes does not require the cross-world joint distribution. Conditioning on post-treatment migration, survival, enrollment or employment changes the estimand and needs separate justification.

\textbf{Statistical statements.} An estimator, confidence set, forecast or test specifies a sampling law, model class and loss/coverage criterion. Uniform \((1-\alpha)\)-coverage means \(\inf_{m\in\mathcal M}\Pr_m\{\tau(m)\in C_n\}\geq1-\alpha\), or an explicitly asymptotic version. The whole parameter space trivially has coverage, so informativeness requires a stated width, risk or power objective. A forecasting improvement is relative to a named benchmark, information set and evaluation population.

\textbf{Equilibrium and optimization.} An equilibrium comprises feasible plans, optimality, beliefs where needed, budget constraints and all market/resource clearing. The primitives or a stated benchmark define each correspondence \(\mathcal E(p;m)\). Nonemptiness is assumed or proved, never inferred just from compact policy sets. A selected-equilibrium target uses a declared selection rule; a robust target quantifies over all permitted equilibria. Use suprema or infima unless attainment is established. An \(\varepsilon\)-optimal policy is feasible and within \(\varepsilon\) of the finite optimum value. Report an empty feasible set separately.

\textbf{Welfare and accounting.} Normative utility, interpersonal normalization, social weights, numeraire, discounting and the use of public receipts are primitives. Behavioral utility need not coincide with the welfare criterion. Household welfare includes ownership income and rents once; adding profits again to compensating variation generally double-counts them. A monetary equivalent is defined by an explicit transfer and price convention, with monotonicity and range conditions. Average effects, mechanism contributions, transfers and welfare losses are different targets. Infinite sums require convergence and dynamic feasible plans require terminal or no-Ponzi conditions.

\textbf{Source versions.} A working-paper date, revision date and journal publication year are distinguished. A DOI for a journal version is not automatically the DOI of an earlier manuscript. Locators refer to the stated version and printed pages where specified. A metadata-only check does not verify a claim about a full-text conclusion. Every entry's source subsection narrows the provenance claim accordingly.
% END ECONOMICS STATEMENT
\end{document}
